\chapter{Lithium}

\section*{Definition and Overview}
Lithium is a naturally occurring alkali metal and a gold-standard pharmacological agent primarily utilised as a mood stabiliser in the treatment of psychiatric disorders. It is indicated for the acute management and long-term prophylaxis of manic and depressive episodes in bipolar I and II disorders, and as an adjunctive treatment for treatment-resistant major depressive disorder. Despite the advent of newer atypical antipsychotics, lithium remains one of the most effective agents for reducing the risk of suicide in patients with mood disorders. Because of its narrow therapeutic index, lithium therapy requires meticulous clinical monitoring and serum level testing to prevent severe toxicity and manage potential adverse effects.

\section*{Epidemiology}
Lithium has been in clinical use for over seven decades and remains widely prescribed globally. In Australia and other developed countries, bipolar disorder affects approximately 1\% to 2\% of the population. Among these individuals, lithium is prescribed to a significant proportion, with studies indicating that up to 20\% to 30\% of patients with bipolar disorder receive lithium therapy at some point. The drug's usage requires careful patient selection, and it is most commonly initiated in patients aged 18 to 45 years, reflecting the typical age of onset for bipolar disorder. The incidence of lithium-induced chronic kidney disease or hypothyroidism increases with the duration of therapy, particularly in female patients and those aged over 50 years.

\section*{Aetiology and Pathophysiology}
Lithium is a monovalent cation that shares properties with sodium and potassium, allowing it to compete with these ions and alter cellular excitability, signal transduction, and neuroplasticity.
\begin{itemize}
    \item \textbf{Inositol depletion hypothesis:} Lithium inhibits inositol monophosphatase (IMPase), reducing key cellular phosphoinositides, thereby dampening overactive G-protein-coupled receptor signalling pathways associated with mania.
    \item \textbf{Glycogen synthase kinase-3 (GSK-3) inhibition:} Lithium directly inhibits GSK-3, an enzyme involved in cellular transcription, energy metabolism, and neuroplasticity, which promotes neuroprotection and synaptic resilience.
    \item \textbf{Neurotransmitter modulation:} Lithium enhances serotonergic transmission, decreases dopaminergic and glutamatergic activity, and increases the synthesis of GABA, a major inhibitory neurotransmitter, restoring neurotransmitter balance.
    \item \textbf{Neurotrophic effects:} Long-term lithium administration increases the expression of neurotrophic factors, such as brain-derived neurotrophic factor (BDNF), which support neuronal survival and grey matter volume in key brain regions.
    \item \textbf{Renal and thyroid interference:} Lithium competes with sodium in the renal collecting duct, leading to resistance to antidiuretic hormone (ADH), and concentrates in the thyroid gland, where it inhibits thyroid hormone synthesis and release.
\end{itemize}

\section*{Clinical Features}
The clinical features associated with lithium use encompass both its therapeutic benefits and a broad profile of potential adverse effects and signs of toxicity.
\begin{itemize}
    \item \textbf{Therapeutic effects:} Resolution of acute manic symptoms (grandiosity, pressured speech, hyperactivity) within 1 to 2 weeks, and long-term stabilisation of mood with reduced frequency and severity of mood episodes.
    \item \textbf{Common early adverse effects:} Fine hand tremor, mild polyuria and polydipsia, nausea, abdominal discomfort, diarrhoea, and a metallic taste in the mouth, which often improve over time.
    \item \textbf{Endocrine features:} Symptoms of hypothyroidism (fatigue, weight gain, cold intolerance, dry skin) or goitre, occurring in up to 10\% to 20\% of patients on long-term therapy.
    \item \textbf{Renal features:} Nephrogenic diabetes insipidus (characterised by profound polyuria and compensatory polydipsia) and a slow, progressive decline in glomerular filtration rate (GFR) with chronic use.
    \item \textbf{Mild to moderate toxicity features:} Worsening coarse hand tremor, muscle weakness, ataxia, dysarthria, persistent vomiting, watery diarrhoea, and drowsiness (typically corresponding to serum levels of 1.5 to 2.0~mmol/L).
    \item \textbf{Severe toxicity features:} Altered mental status, confusion, seizures, focal neurological deficits, hyperreflexia, cardiac arrhythmias, and coma (typically occurring at serum levels above 2.0~mmol/L, posing an immediate threat to life).
\end{itemize}

\section*{Differential Diagnosis}
When evaluating patients on lithium who present with neurological, renal, or endocrine symptoms, clinicians must differentiate lithium-related adverse effects and toxicity from other medical and psychiatric conditions.
\begin{itemize}
    \item \textbf{Essential tremor or anxiety:} May cause a fine bilateral hand tremor similar to lithium-induced tremor, but is not associated with elevated lithium levels or other systemic toxicity symptoms.
    \item \textbf{Primary hypothyroidism:} Presents with identical clinical and biochemical features (elevated TSH, low free T4) but occurs independently of lithium exposure, requiring confirmation through medication history.
    \item \textbf{Central diabetes insipidus:} Causes polyuria and polydipsia, but responds to desmopressin administration, unlike lithium-induced nephrogenic diabetes insipidus where the kidneys are resistant to ADH.
    \item \textbf{Antipsychotic-induced extrapyramidal symptoms (EPS):} Tremor, rigidity, and akathisia caused by co-prescribed neuroleptics, which must be distinguished from lithium-induced neurotoxicity.
    \item \textbf{Sepsis or acute gastroenteritis:} Can present with confusion, vomiting, diarrhoea, and dehydration, which themselves can precipitate acute lithium toxicity by reducing renal clearance.
\end{itemize}

\section*{When to Seek Medical Care}
Patients taking lithium require immediate medical attention if they experience symptoms of toxicity or severe dehydration, which can rapidly elevate lithium levels to dangerous thresholds.

\textbf{Call 000 (or local emergency services) for:}
\begin{itemize}
    \item Severe confusion, disorientation, or hallucinations.
    \item Seizures, loss of consciousness, or inability to stand or walk due to severe ataxia.
    \item Persistent muscle twitching, severe slurring of speech, or cardiac arrhythmias (palpitations, syncope).
\end{itemize}
Other non-emergency reasons to seek medical care include persistent diarrhoea or vomiting, signs of hypothyroidism, or worsening hand tremor that interferes with daily tasks.

\section*{Diagnosis and Investigations}
The safe administration of lithium relies on structured pre-treatment screening and regular longitudinal monitoring of serum lithium levels, renal function, and thyroid parameters.
\begin{itemize}
    \item \textbf{Serum lithium monitoring:} Regular blood tests drawn exactly 12 hours after the last dose (trough level), aiming for a target therapeutic range of 0.6 to 0.8~mmol/L for maintenance, or 0.8 to 1.0~mmol/L for acute mania.
    \item \textbf{Renal function tests (eGFR and electrolytes):} Measured pre-treatment and every 3 to 6 months to monitor for a decline in renal function, which increases the risk of lithium accumulation.
    \item \textbf{Thyroid function tests (TSH and free T4):} Assessed pre-treatment and every 6 to 12 months to detect lithium-induced hypothyroidism or hyperthyroidism early.
    \item \textbf{Electrocardiogram (ECG):} Performed before starting therapy and annually in patients over 40 years or those with cardiovascular history, to screen for T-wave changes, sinus node dysfunction, or QT prolongation.
    \item \textbf{Serum calcium and parathyroid hormone (PTH):} Monitored periodically to detect lithium-induced hyperparathyroidism and hypercalcaemia.
\end{itemize}

\section*{Management}
The management of lithium therapy involves therapeutic drug monitoring, dose adjustments, and the immediate treatment of toxicity.
\begin{itemize}
    \item \textbf{Dosing and administration:} Typically administered as a single daily dose or in divided doses using immediate or sustained-release formulations, adjusted based on serum trough levels.
    \item \textbf{Hydration and sodium management:} Ensuring a stable intake of water and dietary sodium, as dehydration, low-sodium diets, and medications like NSAIDs, ACE inhibitors, or diuretics can reduce lithium excretion and cause toxicity.
    \item \textbf{Managing mild toxicity:} Stop lithium immediately, check serum levels, administer intravenous normal saline (0.9\% NaCl) to promote renal lithium clearance, and monitor electrolytes.
    \item \textbf{Haemodialysis:} The treatment of choice for severe lithium toxicity (typically levels above 4.0~mmol/L, or above 2.5~mmol/L with severe neurological impairment or renal failure) to rapidly remove lithium from the circulation.
    \item \textbf{Treating endocrine side effects:} Lithium-induced hypothyroidism is treated with levothyroxine replacement therapy, allowing lithium to be continued if it is highly effective for the patient's mood stability.
\end{itemize}

\section*{Complications}
\begin{itemize}
    \item \textbf{Chronic kidney disease (CKD):} Long-term tubulointerstitial nephritis leading to a progressive loss of renal concentrating ability and a decline in GFR, potentially culminating in end-stage kidney disease.
    \item \textbf{Nephrogenic diabetes insipidus (NDI):} Reduced ability of the kidneys to concentrate urine, leading to chronic polyuria and polydipsia, which increases the risk of rapid dehydration.
    \item \textbf{Hypothyroidism and goitre:} Impaired thyroid hormone release leading to clinical hypothyroidism and compensatory enlargement of the thyroid gland.
    \item \textbf{Hyperparathyroidism and hypercalcaemia:} Lithium-induced stimulation of the parathyroid glands, leading to elevated calcium levels and potential bone and renal complications.
    \item \textbf{Irreversible neurotoxicity:} Also known as Syndrome of Irreversible Lithium-Effectuated Neurotoxicity (SILENT), presenting with persistent cerebellar dysfunction and cognitive deficits after severe toxicity.
    \item \textbf{Teratogenicity (Ebstein's anomaly):} A rare congenital heart defect in the fetus when lithium is taken during the first trimester of pregnancy.
\end{itemize}

\section*{Prognosis and Outcomes}
When properly monitored, lithium has an excellent therapeutic profile, with up to 60\% to 70\% of patients experiencing a positive clinical response and a dramatic reduction in manic and depressive relapses. The long-term prognosis is highly favourable for patients who adhere to treatment and attend scheduled blood monitoring. However, the risk of renal insufficiency and endocrine dysfunction requires lifelong vigilance, and prompt recognition of toxicity is essential to prevent permanent neurological or cardiovascular damage.

\section*{Prevention and Self-Care}
\begin{itemize}
    \item \textbf{Maintain fluid intake:} Drink 2 to 3 litres of water daily, especially during exercise, hot weather, or illness, to prevent dehydration.
    \item \textbf{Keep sodium intake consistent:} Avoid sudden changes in dietary salt intake, such as starting a low-sodium diet, which can decrease lithium excretion.
    \item \textbf{Avoid certain medications:} Do not take non-steroidal anti-inflammatory drugs (NSAIDs) like ibuprofen, or diuretics, without consulting your doctor, as they can rapidly increase lithium levels.
    \item \textbf{Recognise early toxicity symptoms:} Be vigilant for signs like nausea, vomiting, diarrhoea, muscle weakness, or a coarse hand tremor, and seek immediate medical advice if they occur.
    \item \textbf{Adhere to blood tests:} Attend all scheduled laboratory appointments for serum lithium, kidney, and thyroid tests.
    \item \textbf{Carry medical identification:} Wear a medical alert bracelet or carry a card indicating that you are taking lithium.
\end{itemize}

\section*{Sources and Further Reading}
The following reputable organisations provide further information and clinical guidelines on lithium therapy and its monitoring:
\begin{itemize}
    \item Royal Australian and New Zealand College of Psychiatrists (RANZCP). \textit{Clinical practice guidelines for bipolar disorder.} https://www.ranzcp.org/
    \item National Institute for Health and Care Excellence (NICE). \textit{Bipolar disorder: assessment and management.} https://www.nice.org.uk/
    \item NPS MedicineWise. \textit{Lithium therapy: what you need to know.} https://www.nps.org.au/
    \item Mayo Clinic. \textit{Lithium (oral route) proper use.} https://www.mayoclinic.org/
    \item Head to Health (Australian Department of Health). \textit{Bipolar disorder treatment options.} https://www.headtohealth.gov.au/
\end{itemize}
